\section{Computer-generated samples}

Computer-generated samples are useful because they let us study data when the
underlying probability model is known. In real data analysis, the true model is
usually unknown. In simulation, we choose the model ourselves and then observe
what finite samples from that model look like.

This makes simulation an important learning tool. It helps us see the difference
between a theoretical distribution and the empirical summaries computed from one
finite sample.

\subsection*{What simulation does}

Suppose we ask a computer to generate
\[
X_1,X_2,\ldots,X_n
\]
from a standard normal distribution. The model has
\[
\E[X]=0,
\qquad
\Var(X)=1.
\]
After generation, we obtain observed values
\[
x_1,x_2,\ldots,x_n.
\]
From these values we compute
\[
\bar{x},
\qquad
v_n,
\qquad
F_n,
\]
and draw a histogram.

The sample mean will usually be close to \(0\) for large \(n\), but it will not
usually be exactly \(0\). The histogram may resemble the normal density, but it
will not be identical. The empirical CDF may resemble the normal CDF, but it will
be a step function.

\subsection*{A basic Python-style workflow}

A typical computational experiment has the following structure:
\begin{verbatim}
import numpy as np

sample = np.random.normal(loc=0, scale=1, size=1000)
sample_mean = np.mean(sample)
sample_std = np.std(sample)
q1, median, q3 = np.quantile(sample, [0.25, 0.5, 0.75])
\end{verbatim}

The commands are not the main point. The mathematical point is that the computer
creates one finite sample from a specified model, and then the empirical
summaries describe that sample.

\subsection*{Changing the sample size}

One of the most useful experiments is to compare samples of different sizes. For
example, generate samples from the same distribution with
\[
n=20,\qquad n=100,\qquad n=1000.
\]
For each sample, compute the sample mean, draw a histogram, and plot the
empirical CDF.

Typically, small samples look irregular. Larger samples often display the shape
of the underlying distribution more clearly. The sample mean often gets closer to
the theoretical expectation, and the empirical CDF often gets closer to the
theoretical CDF.

This is an observed pattern, not yet a proof. Later chapters will explain it
using the law of large numbers and related limit results.

\subsection*{Comparing several models}

Simulation also helps compare different probability models. For example:
\begin{itemize}
    \item samples from a uniform distribution tend to fill an interval evenly;
    \item samples from a normal distribution tend to cluster near the center and
    produce a bell-shaped histogram;
    \item samples from an exponential distribution are non-negative and often
    right-skewed;
    \item samples from a beta distribution lie between \(0\) and \(1\) and can
    take many shapes depending on the parameters.
\end{itemize}

Seeing these samples makes the earlier theoretical distributions less abstract.
The density describes the ideal shape. The simulated sample shows one finite
realization from that ideal model.

\subsection*{Simulation is not proof}

A computer experiment can suggest a pattern, illustrate a definition, or reveal a
mistake in intuition. It does not replace mathematical proof. For instance, if
we generate many samples and see their means close to \(0\), this supports the
idea that sample means tend to stabilize, but it does not prove the law of large
numbers.

Simulation is best used together with definitions and theory:
\begin{itemize}
    \item definitions tell us what to compute;
    \item simulation shows how the computations behave in examples;
    \item theory explains why the behaviour occurs under assumptions.
\end{itemize}

\begin{summarybox}
Computer-generated samples help connect probability models with empirical data.
They show how histograms, empirical CDFs, sample means, variances, and quantiles
behave for finite samples from known distributions. Simulation illustrates
mathematics; it does not replace mathematical justification.
\end{summarybox}
